\documentclass[11pt]{article}
\usepackage{mathblatt}
\usepackage{multicol}


\begin{document}
\pfheftstil
\pfheftkopf{Dreiecke \textperiodcentered{} Lösungen}{}
\begin{pfloesung}{}
\lz{1.}{c (liegt dem rechten Winkel gegenüber)}{}{}
\end{pfloesung}
\begin{pfloesung}{}
\lz{2.}{100; 10; 144; 12}{}{}
\end{pfloesung}
\begin{pfloesung}{}
\lz{3.}{ja; ja; nein}{}{}
\end{pfloesung}
\begin{pfloesung}{}
\lz{4.}{Gegenkathete a; Ankathete b; Hypotenuse c}{}{}
\end{pfloesung}
\begin{pfloesung}{}
\lz{5.}{sin; cos; tan}{}{}
\end{pfloesung}
\begin{pfloesung}{}
\lz{6.}{x = 4; x = 5; x = 12}{}{}
\end{pfloesung}
\begin{pfloesung}{}
\lz{7.}{$\gamma$ = 70°; Seite c}{}{}
\end{pfloesung}
\begin{pfloesung}{}
\lz{8.}{65°; 106°; 53,7°}{}{}
\end{pfloesung}
\begin{pfloesung}{}
\lz{9.}{ja; nein}{}{}
\end{pfloesung}
\begin{pfloesung}{}
\lz{10.}{Hypotenuse: $41$ cm}{}{}
\end{pfloesung}
\begin{pfloesung}{}
\lz{11.}{plus}{Hypotenuse gesucht}{}
\end{pfloesung}
\begin{pfloesung}{}
\lz{12.}{Die Summe der Flächeninhalte der Kathetenquadrate ist gleich dem Flächeninhalt des Hypotenusenquadrates}{a² + b² = c²}{}
\end{pfloesung}
\begin{pfloesung}{}
\lz{13.}{z² = x² + y²}{}{}
\end{pfloesung}
\begin{pfloesung}{}
\lz{14.}{z = √(x² + y²)}{}{}
\end{pfloesung}
\begin{pfloesung}{}
\lz{15.}{z = √(x² $-$ y²) (vierte Option)}{z² = x² $-$ y²}{}
\end{pfloesung}
\begin{pfloesung}{}
\lz{16.}{w = √(u² + v²) (dritte Option)}{w² = u² + v²}{}
\end{pfloesung}
\begin{pfloesung}{}
\lz{17.}{$8{,}5$ cm}{}{}
\end{pfloesung}
\begin{pfloesung}{}
\lz{18.}{Katheten: die Höhe und der Überstand}{}{}
\end{pfloesung}
\begin{pfloesung}{}
\lz{19.}{AB = 10√13 $\approx$ 36,1 m}{AB² = 12² + 34² = 1300}{}
\end{pfloesung}
\begin{pfloesung}{}
\lz{20.}{AD = √553 719 $\approx$ 744 m}{AD² = 920² $-$ 541²}{}
\end{pfloesung}
\begin{pfloesung}{}
\lz{21.}{BC = √65 $\approx$ 8,06 m}{BC² = 9² $-$ 4² = 65}{}
\end{pfloesung}
\begin{pfloesung}{}
\lz{22.}{AD = √(14,1² $-$ 7,4²) $\approx$ 12,0 m}{198,81 $-$ 54,76 = 144,05}{}
\end{pfloesung}
\begin{pfloesung}{}
\lz{23.}{FA = 3√9159 $\approx$ 287,1 m}{FA² = 384² $-$ 255² = 82 431}{}
\end{pfloesung}
\begin{pfloesung}{}
\lz{24.}{√89,96 + 2 $\approx$ 11,5 cm}{Diagonale √(7² + 6,4²) = √89,96 $\approx$ 9,5 cm; Durchmesser 6,4 cm}{}
\end{pfloesung}
\begin{pfloesung}{}
\lz{25.}{$9$ cm}{Radius}{}
\end{pfloesung}
\begin{pfloesung}{}
\lz{26.}{rechter Winkel bei C}{ja}{}
\end{pfloesung}
\begin{pfloesung}{}
\lz{27.}{Radius 3,2 m und Dachhöhe als Katheten, Mantellinie als Hypotenuse $\Rightarrow$ s = √(3,2² + h²)}{}{}
\end{pfloesung}
\begin{pfloesung}{}
\lz{28.}{25 + √51,87 $\approx$ 32,2 m}{Kegelhöhe = √(8,6² $-$ 4,7²) = √51,87 $\approx$ 7,20 m}{}
\end{pfloesung}
\begin{pfloesung}{}
\lz{29.}{x = 2√7289 $\approx$ 170,8 cm; $\alpha$ = tan$^{-1}$($\tfrac{8}{85}$) $\approx$ 5,4°}{x² = 170² + 16² = 29 156; tan $\alpha$ = 16 : 170 = $\tfrac{8}{85}$ $\approx$ 0,0941}{}
\end{pfloesung}
\begin{pfloesung}{}
\lz{30.}{BC = √20 $\approx$ 4,47; $\beta$ = tan$^{-1}$(2) $\approx$ 63,4°}{AB = 2, AC = 4; BC² = 2² + 4² = 20; tan $\beta$ = 4 : 2 = 2}{}
\end{pfloesung}
\begin{pfloesung}{}
\lz{31.}{im Zähler: mal}{}{}
\end{pfloesung}
\begin{pfloesung}{}
\lz{32.}{r: G, s: A, t: H}{}{}
\end{pfloesung}
\begin{pfloesung}{}
\lz{33.}{x = 14}{}{}
\end{pfloesung}
\begin{pfloesung}{}
\lz{34.}{sin $\beta$ = r/t}{}{}
\end{pfloesung}
\begin{pfloesung}{}
\lz{35.}{sin $\alpha$ = r/t}{}{}
\end{pfloesung}
\begin{pfloesung}{}
\lz{36.}{tan $\alpha$ = r/s}{}{}
\end{pfloesung}
\begin{pfloesung}{}
\lz{37.}{sin $\gamma$ = u/w}{sin = Gegenkathete : Hypotenuse}{}
\end{pfloesung}
\begin{pfloesung}{}
\lz{38.}{sin $\beta$ = b/a}{}{}
\end{pfloesung}
\begin{pfloesung}{}
\lz{39.}{Winkel: Umkehrtaste}{}{}
\end{pfloesung}
\begin{pfloesung}{}
\lz{40.}{sin}{}{}
\end{pfloesung}
\begin{pfloesung}{}
\lz{41.}{cos $\alpha$ = $\tfrac{4}{9}$ $\Rightarrow$ $\alpha$ $\approx$ 63,6° $\approx$ 64°}{cos $\alpha$ = $\tfrac{4}{9}$ $\approx$ 0,444}{}
\end{pfloesung}
\begin{pfloesung}{}
\lz{42.}{$\alpha$ = sin$^{-1}$($\tfrac{74}{141}$) $\approx$ 31,7°}{sin $\alpha$ = 7,4 : 14,1 = $\tfrac{74}{141}$ $\approx$ 0,525}{}
\end{pfloesung}
\begin{pfloesung}{}
\lz{43.}{$\beta$1 = cos$^{-1}$($\tfrac{85}{128}$) $\approx$ 48,4°}{cos $\beta$1 = 255 : 384 = $\tfrac{85}{128}$ $\approx$ 0,664}{}
\end{pfloesung}
\begin{pfloesung}{}
\lz{44.}{$\alpha$ = tan$^{-1}$($\tfrac{14}{9}$) $\approx$ 57,3°}{15 cm $-$ 6 cm = 9 cm; tan $\alpha$ = 14 : 9}{}
\end{pfloesung}
\begin{pfloesung}{}
\lz{45.}{Teildreieck aus Seite $b$, Höhe und Stück von $a$: $b$ H, Höhe G, Stück von $a$ A}{}{}
\end{pfloesung}
\begin{pfloesung}{}
\lz{46.}{linkes Teildreieck: $b$ H, Höhe G, linkes Stück von $c$ A}{}{}
\end{pfloesung}
\begin{pfloesung}{}
\lz{47.}{57,9 m; 69,9 m}{CB = 112,8 · sin 30° = 56,4 m; 56,4 m + 1,5 m = 57,9 m}{}
\end{pfloesung}
\begin{pfloesung}{}
\lz{48.}{BD = 4,1 · sin 123° : sin 36° $\approx$ 5,85 cm}{BD : sin 123° = 4,1 : sin 36° $\Rightarrow$ BD $\approx$ 5,85 cm}{}
\end{pfloesung}
\begin{pfloesung}{}
\lz{49.}{PB = 1,20 : tan 34,9° $\approx$ 1,72 m}{tan 34,9° = 1,20 : PB}{}
\end{pfloesung}
\begin{pfloesung}{}
\lz{50.}{AC = 12 · sin 55° : sin 5° $\approx$ 112,8 m}{AC : sin 55° = 12 : sin 5°}{}
\end{pfloesung}
\begin{pfloesung}{}
\lz{51.}{AB = 1500 · sin 55° $\approx$ 1228,7 m $\approx$ 1229 m}{sin 55° = AB : 1500}{}
\end{pfloesung}
\begin{pfloesung}{}
\lz{52.}{AD = 1500 · cos 55° $\approx$ 860,4 m}{cos 55° = AD : 1500}{}
\end{pfloesung}
\begin{pfloesung}{}
\lz{53.}{a = 7,4 : sin 52° $\approx$ 9,4 m}{sin 52° $\approx$ 0,788}{}
\end{pfloesung}
\begin{pfloesung}{}
\lz{54.}{y = 160 · sin 141° : sin 34° $\approx$ 180,1 cm}{dritter Winkel 34° gegenüber 160 cm}{}
\end{pfloesung}
\begin{pfloesung}{}
\lz{55.}{2800 · sin 60° : sin 50° + 2500 $\approx$ 5665 m ($\approx$ 5,67 km)}{JR = 2800 · sin 60° : sin 50° $\approx$ 3165 m}{}
\end{pfloesung}
\begin{pfloesung}{}
\lz{56.}{BD = √553 719 · tan 34° $\approx$ 502 m}{$\sphericalangle$CAD = 36°; $\sphericalangle$DAB = 70° $-$ 36° = 34°; BD = AD · tan 34°}{}
\end{pfloesung}
\begin{pfloesung}{}
\lz{57.}{BC = 384 · sin 38° : sin 34° $\approx$ 422,8 m}{$\gamma$ = 180° $-$ 38° $-$ 108° = 34°}{}
\end{pfloesung}
\begin{pfloesung}{}
\lz{58.}{BF = 131,5 · sin 45° = 65,75√2 $\approx$ 93,0 cm = AF; BC = 65,75√2 : cos 65° $\approx$ 220,0 cm}{sin 45° = √$\tfrac{2}{2}$ $\approx$ 0,7071; cos 65° $\approx$ 0,4226}{}
\end{pfloesung}
\begin{pfloesung}{}
\lz{59.}{DP = 11 · sin 115° : sin 49° $-$ 10 $\approx$ 3,2 m}{$\sphericalangle$D = 180° $-$ 115° $-$ 16° = 49°; DE : sin 115° = 11 : sin 49° $\Rightarrow$ DE $\approx$ 13,2 m}{}
\end{pfloesung}
\begin{pfloesung}{}
\lz{60.}{BC = 2004 · sin 60° : sin 74° $\approx$ 1805,5 m; Aussage falsch (DC = 2004 m)}{$\sphericalangle$DBC = 109° $-$ 35° = 74°; BC : sin 60° = 2004 : sin 74°}{}
\end{pfloesung}
\begin{pfloesung}{}
\lz{61.}{AD = 9 · sin 60° : sin 80° $\approx$ 7,9 m}{$\sphericalangle$ACB = 90° $-$ 64° = 26°; $\sphericalangle$ACD = 86° $-$ 26° = 60°; AD : sin 60° = 9 : sin 80°}{}
\end{pfloesung}
\begin{pfloesung}{}
\lz{62.}{BC = 4,5 · sin 62,8° : sin 34,9° $\approx$ 7,00 m; 8 · (4,5 + BC) $\approx$ 92 m²}{BC : sin 62,8° = 4,5 : sin 34,9° $\Rightarrow$ BC $\approx$ 7,00 m; Rechtecke 36 m² und 56 m²}{}
\end{pfloesung}
\begin{pfloesung}{}
\lz{63.}{Sinussatz}{}{}
\end{pfloesung}
\begin{pfloesung}{}
\lz{64.}{57,9 m; 69,9 m}{CB = 112,8 · sin 30° = 56,4 m; 56,4 m + 1,5 m = 57,9 m}{}
\end{pfloesung}
\begin{pfloesung}{}
\lz{65.}{BD = 4,1 · sin 123° : sin 36° $\approx$ 5,85 cm}{BD : sin 123° = 4,1 : sin 36° $\Rightarrow$ BD $\approx$ 5,85 cm}{}
\end{pfloesung}
\begin{pfloesung}{}
\lz{66.}{PB = 1,20 : tan 34,9° $\approx$ 1,72 m}{tan 34,9° = 1,20 : PB}{}
\end{pfloesung}
\begin{pfloesung}{}
\lz{67.}{AC = 12 · sin 55° : sin 5° $\approx$ 112,8 m}{AC : sin 55° = 12 : sin 5°}{}
\end{pfloesung}
\begin{pfloesung}{}
\lz{68.}{AB = 1500 · sin 55° $\approx$ 1228,7 m $\approx$ 1229 m}{sin 55° = AB : 1500}{}
\end{pfloesung}
\begin{pfloesung}{}
\lz{69.}{AD = 1500 · cos 55° $\approx$ 860,4 m}{cos 55° = AD : 1500}{}
\end{pfloesung}
\begin{pfloesung}{}
\lz{70.}{a = 7,4 : sin 52° $\approx$ 9,4 m}{sin 52° $\approx$ 0,788}{}
\end{pfloesung}
\begin{pfloesung}{}
\lz{71.}{y = 160 · sin 141° : sin 34° $\approx$ 180,1 cm}{dritter Winkel 34° gegenüber 160 cm}{}
\end{pfloesung}
\begin{pfloesung}{}
\lz{72.}{2800 · sin 60° : sin 50° + 2500 $\approx$ 5665 m ($\approx$ 5,67 km)}{JR = 2800 · sin 60° : sin 50° $\approx$ 3165 m}{}
\end{pfloesung}
\begin{pfloesung}{}
\lz{73.}{BD = √553 719 · tan 34° $\approx$ 502 m}{$\sphericalangle$CAD = 36°; $\sphericalangle$DAB = 70° $-$ 36° = 34°; BD = AD · tan 34°}{}
\end{pfloesung}
\begin{pfloesung}{}
\lz{74.}{BC = 384 · sin 38° : sin 34° $\approx$ 422,8 m}{$\gamma$ = 180° $-$ 38° $-$ 108° = 34°}{}
\end{pfloesung}
\begin{pfloesung}{}
\lz{75.}{BF = 131,5 · sin 45° = 65,75√2 $\approx$ 93,0 cm = AF; BC = 65,75√2 : cos 65° $\approx$ 220,0 cm}{sin 45° = √$\tfrac{2}{2}$ $\approx$ 0,7071; cos 65° $\approx$ 0,4226}{}
\end{pfloesung}
\begin{pfloesung}{}
\lz{76.}{DP = 11 · sin 115° : sin 49° $-$ 10 $\approx$ 3,2 m}{$\sphericalangle$D = 180° $-$ 115° $-$ 16° = 49°; DE : sin 115° = 11 : sin 49° $\Rightarrow$ DE $\approx$ 13,2 m}{}
\end{pfloesung}
\begin{pfloesung}{}
\lz{77.}{BC = 2004 · sin 60° : sin 74° $\approx$ 1805,5 m; Aussage falsch (DC = 2004 m)}{$\sphericalangle$DBC = 109° $-$ 35° = 74°; BC : sin 60° = 2004 : sin 74°}{}
\end{pfloesung}
\begin{pfloesung}{}
\lz{78.}{AD = 9 · sin 60° : sin 80° $\approx$ 7,9 m}{$\sphericalangle$ACB = 90° $-$ 64° = 26°; $\sphericalangle$ACD = 86° $-$ 26° = 60°; AD : sin 60° = 9 : sin 80°}{}
\end{pfloesung}
\begin{pfloesung}{}
\lz{79.}{BC = 4,5 · sin 62,8° : sin 34,9° $\approx$ 7,00 m; 8 · (4,5 + BC) $\approx$ 92 m²}{BC : sin 62,8° = 4,5 : sin 34,9° $\Rightarrow$ BC $\approx$ 7,00 m; Rechtecke 36 m² und 56 m²}{}
\end{pfloesung}
\begin{pfloesung}{}
\lz{80.}{1~Pythagoras, 2~Winkelfunktion, 3~Seite berechnen, 4~Winkel, 5~Pythagoras, 6~Winkelfunktion}{}{}
\end{pfloesung}
\begin{pfloesung}{}
\lz{81.}{gleichschenklig}{}{}
\end{pfloesung}
\begin{pfloesung}{}
\lz{82.}{Scheitelwinkel}{}{}
\end{pfloesung}
\begin{pfloesung}{}
\lz{83.}{$\beta$ = 105°}{180° $-$ 75°}{}
\end{pfloesung}
\begin{pfloesung}{}
\lz{84.}{$\alpha$ = 40°}{}{}
\end{pfloesung}
\begin{pfloesung}{}
\lz{85.}{$\gamma$ = 82,3°}{}{}
\end{pfloesung}
\begin{pfloesung}{}
\lz{86.}{$\delta$ = 21°}{}{}
\end{pfloesung}
\begin{pfloesung}{}
\lz{87.}{$\gamma$ = 124°}{180° $-$ 56°}{}
\end{pfloesung}
\begin{pfloesung}{}
\lz{88.}{$\sphericalangle$ACB = 60° $\Rightarrow$ $\sphericalangle$ACD = 120° $\Rightarrow$ $\delta$ = 180° $-$ 120° $-$ 5° = 55°}{Winkel ACD = 120°}{}
\end{pfloesung}
\begin{pfloesung}{}
\lz{89.}{$\gamma_1$ = cos$^{-1}$($\tfrac{74}{141}$) $\approx$ 58,3°}{180° $-$ 90° $-$ 31,7°}{}
\end{pfloesung}
\begin{pfloesung}{}
\lz{90.}{$2$ Achsen: die beiden Mittellinien}{}{}
\end{pfloesung}
\begin{pfloesung}{}
\lz{91.}{2}{}{}
\end{pfloesung}
\begin{pfloesung}{}
\lz{92.}{4}{}{}
\end{pfloesung}
\begin{pfloesung}{}
\lz{93.}{1}{Mittelsenkrechte der parallelen Seiten}{}
\end{pfloesung}
\begin{pfloesung}{}
\lz{94.}{1 Symmetrieachse (BD); BF = 80 $-$ 25 = 55 cm; AB = 5√146 $\approx$ 60,4 cm}{AF = 25 cm; AB² = 25² + 55² = 3650 cm²}{}
\end{pfloesung}
\begin{pfloesung}{}
\lz{95.}{Umkehrung: Wenn ein Tier vier Beine hat, dann ist es ein Hund}{}{}
\end{pfloesung}
\begin{pfloesung}{}
\lz{96.}{$40$ cm}{Hypotenuse ist $KM$; $KM = 2{,}6$ cm; Kathete gesucht: $KL^2 = 42{,}25 - 10{,}89$, $KL = 5{,}6$ dm; $KM = 41$ cm, $LM^2 = 1\,681 - 81$, $LM$; gegenüber $L$; a; b}{}
\end{pfloesung}
\begin{pfloesung}{}
\lz{97.}{$\approx 46{,}6^\circ$}{Sinus, $\alpha \approx 25{,}3^\circ$; Tangens mit $60$ cm, $\alpha \approx 36{,}9^\circ$; Kosinus mit $16$ cm, $\alpha$; 1; 2; 3}{}
\end{pfloesung}
\begin{pfloesung}{}
\lz{98.}{$24$ cm}{$d = \sqrt{36 + 4 + 81} = 11$ cm; Überstand $5$ cm, $h = \sqrt{13^2 - 5^2} = 12$ cm; Radius $7$ cm, $h = \sqrt{25^2 - 7^2}$; 1; 3}{}
\end{pfloesung}
\begin{pfloesung}{}
\lz{99.}{$\approx 6{,}2$ cm}{rechtwinklig, Sinus: $a = 8 \cdot \mathrm{sin}\,40^\circ$ $\approx 5{,}1$ cm; Sinussatz: $b = \tfrac{6 \cdot \mathrm{sin}\,75^\circ}{\mathrm{sin}\,50^\circ}$ $\approx 7{,}6$ cm; Kosinussatz: $c = \sqrt{39}$ ; 1; 3}{}
\end{pfloesung}
\begin{pfloesung}{}
\lz{100.}{Dreieck ADC}{}{}
\end{pfloesung}
\begin{pfloesung}{}
\lz{101.}{gegenüberliegende Winkel gleich groß}{}{}
\end{pfloesung}
\begin{pfloesung}{}
\lz{102.}{gibt es ein Paar paralleler Seiten}{Definition Trapez}{}
\end{pfloesung}
\begin{pfloesung}{}
\lz{103.}{BC = 6 cm ($\alpha$ = $\beta$ $\Rightarrow$ gleichschenklig); $\gamma$ = 40°}{180° $-$ 2 · 70°}{}
\end{pfloesung}
\begin{pfloesung}{}
\lz{104.}{Winkel bei A und B im Dreieck ABF je 45° $\Rightarrow$ gleichschenklig mit Basis AB (BF = AF)}{Winkelsumme: 180° $-$ 90° $-$ 45° = 45°}{}
\end{pfloesung}
\begin{pfloesung}{}
\lz{105.}{Seitenhalbierende}{}{}
\end{pfloesung}
\begin{pfloesung}{}
\lz{106.}{AF = DF = 25 cm $\Rightarrow$ Dreieck AFD gleichschenklig-rechtwinklig $\Rightarrow$ $\angle$ADF = 45°, ebenso $\angle$FDC = 45° $\Rightarrow$ $\gamma$ = 90° (oder Thales: D auf Kreis um F mit r = 25 cm)}{AF = FC = DF = 25 cm; Basiswinkel 45°}{}
\end{pfloesung}
\begin{pfloesung}{\textbf{Prüfstein}}
\lz{a)}{h = 3√95 $\approx$ 29,2 cm}{h² = 32² $-$ 13² = 855}{}
\lz{b)}{$\varepsilon$ = cos$^{-1}$($\tfrac{13}{32}$) $\approx$ 66,0°}{cos $\varepsilon$ = 13 : 32 = $\tfrac{13}{32}$ $\approx$ 0,406}{}
\lz{c)}{AC = 3√95 : sin 22° $\approx$ 78,1 cm}{sin 22° = h : AC; AF = 3√95 : tan 22° $\approx$ 72,4 cm}{}
\end{pfloesung}
\end{document}